% TOP_PROBLEM: {"id": 3263, "title": "Caccetta-Häggkvist Conjecture", "category_id": 3, "category": {"id": 3, "name": "graph_theory", "display_name": "Graph Theory", "description": "Problems involving graphs, networks, and their properties.", "slug": "graph-theory", "order_index": 3, "created_at": "2026-07-31T15:26:25.671Z"}, "status": "open", "problem_number": "OPG-46385", "set_id": 12}
% FIRST_POSTED: 2026-10-01
% ENTRY_KIND: statement_only
% UnsolvedMath ID 3263; source catalogue code OPG-46385
% !TeX program = lualatex
% Definition and sources only; this file is not an LLM solution attempt.
\documentclass[11pt,a4paper]{article}
\usepackage[margin=25mm]{geometry}
\usepackage{fontspec}
\setmainfont{lmroman10-regular.otf}[BoldFont=lmroman10-bold.otf,
 ItalicFont=lmroman10-italic.otf,BoldItalicFont=lmroman10-bolditalic.otf]
\usepackage{amsmath,amssymb,mathtools}
\usepackage{unicode-math}
\setmathfont{latinmodern-math.otf}
\AtBeginDocument{\let\setminus\smallsetminus}
\usepackage{xurl,hyperref}
\hypersetup{colorlinks=true,urlcolor=blue}
\setcounter{secnumdepth}{0}
\setlength{\parindent}{0pt}
\setlength{\parskip}{5pt}
\setlength{\emergencystretch}{3em}
\begin{document}
\section*{Caccetta-Häggkvist Conjecture}\label{3263}
{\small UnsolvedMath ID 3263 · OPG-46385 · Graph Theory · OpenGarden}
\subsection{Definitions and mathematical statement}
Let $n\ge2$ and $1\le r\le n-1$ be integers. A simple digraph on $n$
vertices has no loops and has at most one arc for each ordered pair of
distinct vertices; opposite arcs are permitted. Its minimum outdegree
$\delta^+(D)$ is the smallest number of distinct outneighbours of a
vertex. A directed cycle of length $\ell\ge2$ has $\ell$ distinct
vertices and follows the arc directions around them; opposite arcs
therefore constitute a cycle of length two.

The Caccetta--Häggkvist conjecture is
\[
 \delta^+(D)\ge r\quad\Longrightarrow\quad
 D\text{ has a directed cycle of length at most }
                         \left\lceil\frac nr\right\rceil.
\]
The ceiling is the least integer at least $n/r$. The quantifiers cover
every finite simple digraph satisfying the degree hypothesis. No
indegree bound, regularity, or connectedness assumption is imposed.
Because every vertex has an outgoing arc, following arcs eventually
produces a cycle; the problem is the specific length guarantee.

For example, when $r\ge n/3$ the assertion asks for a cycle of length
two or three. When the digraph is an orientation and hence has no
2-cycle, this becomes a directed-triangle assertion. The parameter $n$
counts vertices of the whole digraph and $r$ is a uniform lower bound,
not an average degree. Replacing a cycle by a closed walk would weaken
the formulation, although any shortest directed closed walk contains
a simple directed cycle.
\subsection{Short English statement}
Does an n-vertex digraph of minimum outdegree r always contain a directed cycle of length at most the ceiling of n/r?
\subsection{Sources}
\begin{itemize}
\item[S1] ULAM.AI and the UnsolvedMath Contributors, supplied problems.json, record 3263 (OPG-46385), OpenGarden collection. Catalogue identity and source transcription; CC BY 4.0.
\url{https://huggingface.co/datasets/ulamai/UnsolvedMath}.
\item[S2] Open Problem Garden, Caccetta-Häggkvist Conjecture. Original problem and discussion; the mathematical formulation below is expanded from the supplied source record.
\url{http://www.openproblemgarden.org/op/caccetta_haggkvist_conjecture}.
\end{itemize}
\end{document}
